\section{Diskussion}
\label{sec:Diskussion}
Zur Bestimmung der relativen Abweichungen zwischen theoretischem und experimentell bestimmtem Wert wird die Gleichung
\begin{equation*}
    \symup{\Delta}x = \frac{|x_{\mathrm{exp}} - x_{\mathrm{theo}}|}{x_{\mathrm{theo}}}
\end{equation*}
verwendet.
\subsection{invertierender Linearverstärker}
Die theoretischen und expermentellen Werte der Leerlaufverstärkung sowie ihre relative Abweichung sind für die vier Messreihen in \autoref{tab:leer} aufgelistet.
\begin{table}
  \centering
  \caption{Experimentelle und theoretische Werte für die Leerlaufverstärkung, sowie deren relative Abweichung.}
  \label{tab:leer}
  \sisetup{table-format=1.1, per-mode=reciprocal}
  \begin{tblr}{
      colspec = {S[table-format=1.0] S[table-format=2.2] S[table-format=1.3] S[table-format=1.3] S[table-format=1.1] },
      row{1} = {guard, mode=math},
      vline{4} = {2}{-}{text=\clap{$\pm$}},
    }
    \toprule
    \mathrm{Messreihe} & V_{\mathrm{theo}} & \SetCell[c=2]{c} V_{\mathrm{exp}} & & \symup{\Delta}V \mathbin{/} \% \\
    \midrule
    1   & 10    & 9.47 & 0.02   & 5.3  \\
    2   & 6.67  & 6.36 & 0.07   & 4.6  \\
    3   & 2.13  & 1.966 & 0.032 & 6.3  \\
    4   & 1.47  & 1.39 & 0.04   & 5.4  \\
    \bottomrule
  \end{tblr}
\end{table}
Die Abweichungen liegen alle im einstelligen Prozentbereich. Die berechneten Werte stimmen also gut mit der Theorie überein. Ein Zusammenhang zwischen größerer Leerlaufverstärkung und Fehler ist nicht erkennbar.\\
Ein Grund für die Abweichungen kann sein, dass die Plateaus alle eine leichte Steigung aufweisen. Diese ist durch den Fit-Parameter $b_0$ beschrieben. 
In den Berechnungen von \autoref{subsec:linear} ist zu sehen, dass die Parameter $b_0$ für alle vier Messreihen in der Größenordnung $b_0 \sim 10^{-3}$ liegen.
Im Idealfall wäre dieser Wert Null. Eine leichte Abweichung kann über mehrere Dekaden eine Abweichung verursachen.
\\
Die experimentell bestimmten Bandbreitenprodukte sind in \autoref{tab:bandwidth} aufgelistet.
\begin{table}
  \centering
  \caption{EExperimentell bestimmte Werte für das Bandbreitenprodukt.}
  \label{tab:bandwidth}
  \sisetup{table-format=1.1, per-mode=reciprocal}
  \begin{tblr}{
      colspec = {S[table-format=1.0] S[table-format=2.2] S[table-format=1.3] S[table-format=1.3] S[table-format=1.1] },
      row{1} = {guard, mode=math},
      vline{3} = {2}{-}{text=\clap{$\pm$}},
    }
    \toprule
    \mathrm{Messreihe} & \SetCell[c=2]{c} \mathrm{Bandbreitenprodukt} \mathbin{/} \unit{\kilo\hertz} \\
    \midrule
    1 & 200 & 60 \\
    2 & 240 & 90 \\
    3 & 240 & 80 \\
    4 & 240 & 70 \\
    \bottomrule
  \end{tblr}
\end{table}
Hierbei ist auffällig, dass die Mittelwerte der letzten drei Messreihen identisch sind, die Standardabweichung aber mit sinkender Verstärkung fällt.
Das Bandbreitenprodukt der ersten Messreihe mit der höchsten Verstärkung weicht deutlich von den anderen ab. 
Eine Fehlerquelle könnte sein, dass die Bandbreite mit steigender Verstärkung geringer wird und dadurch Abweichungen in der Plateaulänge entstehen.
\subsection{Umkehr-Integrator und invertierender Differenzierer}
Für den Umkehr-Integrator sollte nach \autoref{equ:fint} die Steigung der Gerade $m = -1$ sein, dür en Differenzierer gemäß \autoref{equ:fdiff} $m=1$.
Die Fit-Funktionen liefern für die Steigungen
\begin{align*}
    b_{\mathrm{int,exp}} &= \qty{-0.841(0.005)}{}, & b_{\mathrm{diff,exp}} &= \qty{0.895(0.003)}{}. \\
\end{align*}
Die Abweichungen zu den Theoriewerten betragen
\begin{align*}
    \symup{\Delta}b_{\mathrm{int}} &= \qty{15.9(0.5)}{\%}, & \symup{\Delta}b_{\mathrm{diff}} &= \qty{10.5(0.3)}{\%}. \\
\end{align*}
Die großen Abweichungen deuten auf systematische Fehler hin. Das Ablesen der Werte erfolgte mit den Cursorn des Oszilloskops. 
Es ist möglich, dass beim Ablesen Fehler gemacht wurden. Auch Fehler bei der Auswahl der Widerstände und Kondensatoren sind denkbar. 
Die Werte der Bauteile liegen weit auseinander, was zu Fehlern in der Signalübertragung führen könnte.
\subsection{Schmitt-Trigger}
Sowohl in \autoref{fig:schmittSin} als auch in \autoref{fig:schmittTri} ist die Hysterese des Schmitt-Triggers deutlich zu erkennen.\\
Die experimentellen Bestimmungen liefern Werte für die Schwellspannung von 
\begin{align*}
    U_{\mathrm{Schwell,sin}} &= \pm\qty{1.402(0.0005)}{\volt}, & \overline{U}_{\mathrm{Schwell,tri}} &= \pm\qty{1.4031(0.0045)}{\volt}. \\
\end{align*}
Die Abweichung der beiden Werte zueinander ist erkennbar gering. Zm theoretischen Wert $U_{\mathrm{Schwell,theo}} = \pm\qty{1.36}{\volt}$ betragen die relativen Abweichungen
\begin{align*}
    \symup{\Delta}U_{\mathrm{Schwell,sin}} &= \qty{3.10(0.04)}{\%}, & \symup{\Delta}{U}_{\mathrm{Schwell,tri}} &= \qty{3.17(0.33)}{\%}. \\
\end{align*}
Die Abweichungen zum Theoriewert sind also ebenfalls sehr gering. 
Dass die experimentellen Werte sich näher liegen als dem Theoriewert deutet darauf hin, dass der Wert im Experiment überschätzt wurde. 
Die Werte wurden am Oszilloskop abgelesen und sind möglicherweise leicht verrauscht.